%% kapa: exam=2 round=36 year=2025 period=2 subject=appraisal_theory subject_ko=감정평가이론 seq=1 number=1 points=40
%% source: Q-Net 감정평가사 기출문제 (cst00309) artlSeq=5239325 file=raw/1008/5239325/2025년 제36회 감정평가사 2차 감정평가이론.pdf
%% conversion: kordoc 4.24.51 -> tools/gen2.py (auto); review=proofread-gpt-sol; flags=none
\begin{longquestion}{【문제 1】}
다음 자료를 참고하여 물음에 답하시오. (단, 각 물음은 서로 독립적임) (40점) \par
\begin{framedtext}
1. 농업인인 甲은 거주 목적의 단독주택을 신축하기 위해 A광역시 B군에 소재하는 토지 2필지를 매수하고 소유권이전등기를 한 후, 건축비용 마련을 위해 금융회사에 담보대출을 신청함 \par
2. 매매계약일: 2024. 06. 01., 매매대금: 2억원, 소유권이전등기일: 2025. 06. 01. \par
3. 매매대상토지의 내용은 다음과 같음(매매계약일 기준) \par
{\small\centering
\resizebox{\linewidth}{!}{%
\begin{tabular}{|c|c|c|c|c|}
\hline
\makecell{구분} & \makecell{면적(m²)} & \makecell{지목} & \makecell{용도지역} & \makecell{비고} \\
\hline
\makecell{토지 1} & \makecell{200} & \makecell{과수원} & \makecell{자연녹지\newline{}지역} & \makecell{수령 약 20년생 복숭아나무 15주 소재} \\
\hline
\makecell{토지 2} & \makecell{300} & \makecell{임야} & \makecell{자연녹지\newline{}지역} & \makecell{관상수로 식재한 수고 약 10m의 소나무\newline{}10주 소재} \\
\hline
\end{tabular}}\par}

4. 매매대상토지는 서로 접한 토지로서 자연취락과 인접한 남측하향 완경사지대에 위치함 \par
5. 매도인은 매매대상토지를 하나의 대지로 하여 2023. 12. 01. 단독주택 건축허가를 받았고, 건축공사에 착수하지는 않았음 \par
6. 매매계약의 특약사항으로 매매대상토지에 소재하는 나무도 매매대상물건에 포함되며, 매도인은 건축허가와 관련된 일체의 권리에 대해 매수인에게 무상으로 양도하기로 함 \par
7. 금융회사는 매매대상토지에 대해 감정평가사 乙에게 담보목적의 감정평가를 의뢰함 \par
8. 기준시점(2025. 07. 01.) 현재 甲은 B군수로부터 건축관계자명의변경신고필증을 교부받았음 \par
9. 인근지역에 소재하는 단독주택의 표준적인 획지는 대지면적이 500m² 내외임 \par
\end{framedtext}
\medskip\noindent 물음 1) 乙은 대상물건의 일단지 성립여부를 분석하고 있다. 일단지 성립을 긍정하는 입장에서 그 근거를 논하시오. (15점) \par
\medskip\noindent 물음 2) 乙은 대상물건의 최유효이용분석을 하고 있다. 최유효이용의 판정기준 중 경제적 타당성여부를 분석할 때 검토해야 할 사항에 대해 쓰시오. (15점) \par
\medskip\noindent 물음 3) 乙은 본건 매매사례(거래사례)를 토대로 인근지역의 지가수준을 분석하고 있다. 이 경우 검토해야 할 사항에 대해 쓰시오. (10점) \par
\end{longquestion}
